\section{Strain}
\label{sec:strain}

Strain is computed from the displacement results in a window of its own, \ui{Strain
Post-Processing}, called the strain window in this guide (Figure~\ref{fig:strain-window}). It
opens by itself when a run finishes; \ui{Open Strain Window} in the right sidebar opens it at
any time once there are displacement results (without them, the \ui{LOG} says \emph{Run DIC
first -{}- no displacement results to post-process.}). The window has two tabs: \ui{Strain
Field}, described here, and \ui{Analysis}, which places probes, gauges and areas on the results
and plots them (Section~\ref{sec:analysis}).

\screenshot{fig-strain-window.png}{The strain window, \ui{Strain Field} tab, showing
\ui{εyy} around the two holes of the plate in tension. Left: the toolbar, the canvas with
its colorbar and the frame navigator. Right: \ui{STRAIN PARAMETERS}, the buttons, \ui{FIELD} and
\ui{VISUALIZATION}.}{fig:strain-window}

\subsection{The strain window}

The left part repeats the canvas of the main window: a toolbar with the fit, \ui{100\%}, zoom-in
and zoom-out buttons, the canvas with its colorbar, and a frame navigator with the same buttons and
playback speeds as the main one (Section~\ref{sec:canvas}). The toolbar also says why nothing is
drawn when that happens (\emph{Click Compute Strain to show this field.}, \emph{Frame $k$ was not
computed.}), and shows the position and value under the pointer.

The right column is a scrollable stack of sections; click a header to fold or unfold it.
\ui{STRAIN PARAMETERS}, \ui{FIELD} and \ui{VISUALIZATION} start unfolded, \ui{ADVANCED},
\ui{PHYSICAL UNITS} and \ui{LOG} folded. The buttons \ui{Compute Strain}, \ui{Cancel} and
\ui{Export Results} sit between the parameters and \ui{FIELD}, outside any section, so they stay
in reach when the sections are folded.

The strain window is independent of the main window. It has its own frame, so moving its slider
does not move the main canvas, and its own \ui{Show on}, background, opacity and edge options.
Two things are shared: \ui{PHYSICAL UNITS} (Section~\ref{sec:physical-units}), and the colormap
and range of \ui{U}, \ui{V} and \ui{Magnitude}, which are those of \ui{Disp U}, \ui{Disp V} and
\ui{Magnitude} in the main window.

\begin{tipbox}
The window opens on frame~1, the reference, where every field is zero. Move its slider to a later
frame to see the deformation.
\end{tipbox}

Closing the window keeps the strain result; \ui{Open Strain Window} brings it back as you left
it. Closing it during a strain computation cancels the computation.

\subsection{How strain is computed}

Strain comes from a virtual strain gauge (VSG). Around each point, RAFTcorr takes a square window
of $L_g \times L_g$ pixels, the gauge, and fits a polynomial by least squares to each displacement
component $u$ and $v$ over the valid points in it. The slopes of the fit at the centre are the
displacement gradients $\partial u/\partial x$, $\partial u/\partial y$, $\partial v/\partial x$
and $\partial v/\partial y$. They give the deformation gradient
\[
  \mathbf{F} = \mathbf{I} + \nabla\mathbf{u}, \qquad
  \mathbf{F} = \begin{pmatrix} 1 + \partial u/\partial x & \partial u/\partial y \\
  \partial v/\partial x & 1 + \partial v/\partial y \end{pmatrix},
\]
from which every strain field is computed. The gradients are taken with respect to the frame-1
positions, so strain is Lagrangian like the displacement. RAFTcorr's displacement is dense (one
value per pixel), and the strain is computed at every pixel of the region of interest unless you
ask for a coarser spacing (\ui{Output spacing}).

A point is left without strain when less than a quarter of its gauge holds valid points
(weighted by the gauge's weighting), when the valid points are arranged so that the fit is
ill-conditioned (a thin strip, for example), and on isolated parts of the region that hold too few
points for a fit. A point near an edge also loses its strain to the trim described below.

\subsection{STRAIN PARAMETERS}

\screenshot[0.5\textwidth]{fig-strain-params.png}{\ui{STRAIN PARAMETERS} with the
\ui{ADVANCED} section unfolded below it.}{fig:strain-params}

The section holds four settings (Figure~\ref{fig:strain-params}); the less used ones are under
\ui{ADVANCED}, below it.

\begin{description}[leftmargin=1.2cm,style=nextline]
  \item[Method]
    The polynomial fitted in each gauge. \ui{Plane fitting} (the default) is the standard
    first-order gauge: a plane, so the strain is the mean gradient over the gauge.
    \ui{Quadratic fitting} adds the second-order terms; it follows a strain gradient inside the
    gauge better, at the cost of more noise. Consider it with a large gauge over a strong
    gradient, such as the edge of a hole.
  \item[VSG size]
    The side $L_g$ of the gauge, from 3 to 401\,px; default 41\,px. It is odd, so that the gauge
    has a centre pixel; an even value moves to the next odd one. A larger gauge averages more
    points: smoother strain, less noise, but a lower spatial resolution, and peaks narrower than
    the gauge are flattened. A smaller gauge resolves sharper features with more noise. The width
    of the trimmed band below grows with it.
  \item[Trim low-confidence edges]
    The coefficient $\alpha$ of the edge trim, from 0.00 to 1.00 in steps of 0.05; default 0.70.
    See below.
  \item[Strain type]
    The strain measure: \ui{Infinitesimal}, \ui{Eulerian} or \ui{Green-Lagrange} (the default).
    See below.
\end{description}

\subsubsection*{Trim low-confidence edges}
\label{sec:edge-trim}

Next to an edge, part of a gauge falls outside the valid region and the fit sees points on one
side only. Such a one-sided fit is biased by the strain gradient, which is often largest at the
edge (a hole, a notch), and it is noisier because fewer points enter it. The trim leaves these
points empty rather than reporting a value that looks as reliable as the rest.

An edge is any border of the valid region of the frame: the boundary of the region of interest,
a hole or a crack, an area that became untracked, and the border of the displacement grid. A point
is trimmed when an edge enters the central square of its gauge whose half-width is
$\alpha\,(L_g - 1)/2$, the coefficient times the half-width of the gauge. With the default gauge
of 41\,px and $\alpha = 0.70$, this half-width is $0.70 \times 20$\,px $= 14$\,px, and the band
along a straight edge is 14\,px wide. Put another way, a point next to a straight edge is kept
when the edge cuts off at most about $(1-\alpha)/2$ of its gauge width, 15\,\% at the default.

\begin{itemize}
  \item $\alpha = 0$: no trim; every point with a fit is reported.
  \item $\alpha = 0.70$: the default.
  \item $\alpha = 1$: only points whose whole gauge lies inside the valid region are reported.
\end{itemize}

The trim is computed for every frame on that frame's valid region, so the band follows a crack
that grows or an area that is lost. Trimmed points are empty (NaN) in every strain field and in
exported data files. Below the box, \emph{Trimmed: $N$ points ($p$\%)} gives how many points of
the frame on show the computed strain lost to the trim. It describes the strain on show, computed
with the value of the last \ui{Compute Strain}, not the value now in the box. The empty band can
be filled in pictures with \ui{Fill trimmed edges (display only)} (Section~\ref{sec:strain-viz}).

\begin{warnbox}
If the band is wider than half the width of the region of interest, every point would be
trimmed, and \ui{Compute Strain} stops with \emph{Edge trimming would remove every strain point
\ldots{} Reduce the VSG size or lower `Trim low-confidence edges'.} Do one of the two, then
compute again.
\end{warnbox}

\subsubsection*{Strain type}

All three measures are computed from $\mathbf{F}$:

\begin{description}[leftmargin=1.2cm,style=nextline]
  \item[Infinitesimal]
    $\boldsymbol{\varepsilon} = \tfrac12\left(\nabla\mathbf{u} + \nabla\mathbf{u}^{T}\right)$.
    Valid for small strains and small rotations only. It is not zero under a rigid rotation: a
    rotation by an angle $\theta$ reads as
    $\varepsilon_{xx} = \varepsilon_{yy} = \cos\theta - 1$.
  \item[Eulerian]
    The Euler-Almansi strain
    $\mathbf{e} = \tfrac12\left(\mathbf{I} - (\mathbf{F}\mathbf{F}^{T})^{-1}\right)$, which
    refers to the deformed configuration. Zero under any rigid rotation.
  \item[Green-Lagrange]
    $\mathbf{E} = \tfrac12\left(\mathbf{F}^{T}\mathbf{F} - \mathbf{I}\right)$, which refers to
    the reference configuration. Zero under any rigid rotation. The default.
\end{description}

For small strains and small rotations the three agree. Keep \ui{Green-Lagrange} when the specimen
rotates or strains are large; choose \ui{Infinitesimal} to compare with a small-strain
calculation, \ui{Eulerian} to compare with a measure in the deformed configuration.

\subsection{ADVANCED}

\begin{description}[leftmargin=1.2cm,style=nextline]
  \item[Weighting]
    How the points of a gauge count in the fit: \ui{Uniform} (the default), all alike, or
    \ui{Gaussian}, decreasing with the distance from the centre as
    $\exp\!\left(-(\Delta x^2 + \Delta y^2)/2\sigma^2\right)$, where $\Delta x$ and $\Delta y$ are
    the offsets from the centre. A Gaussian gauge counts the points near its centre more than
    those near its border.
  \item[Gaussian $\sigma$]
    The $\sigma$ of the Gaussian weight, up to 200\,px; 0 shows \ui{Auto (Lg / 4)}, a quarter
    of the gauge size (10.25\,px for 41\,px). Enabled only with \ui{Gaussian} weighting.
  \item[Output spacing]
    Keep the strain every $n$ pixels in $x$ and $y$, from 1 to 64\,px; default 1\,px (every
    pixel). Each kept point is still fitted over its whole gauge. A larger spacing gives a
    smaller strain result: less memory and smaller exported files for long sequences.
  \item[Pre-fit smoothing]
    Gaussian smoothing of the displacement, with $\sigma$ in pixels up to 50\,px, before the
    strain fit; \ui{Off} (0, the default) differentiates the measured displacement as it is.
    Only the strain uses the smoothed displacement; the displacement results do not change. The
    \ui{Smoothing} parameter of the main window (Section~\ref{sec:parameters}) is different: it
    smooths the displacement results themselves. Smoothing widens the effective gauge.
  \item[Never fit across cracks or holes]
    On by default. Each connected part of the tracked region is fitted on its own, so no gauge
    spans a crack, a hole or the gap between two regions. With it off, a gauge that straddles a
    narrow crack fits one surface over both faces, and the opening of the crack shows as a band
    of high strain as wide as the gauge.
\end{description}

\subsection{Compute Strain}

\ui{Compute Strain} computes the strain of every computed frame (frame~1 included) with the
parameters on show, in the background. A progress bar and a status line appear below the
buttons: \ui{Starting\ldots}, then \ui{Strain frame $i$/$n$}, which counts the frames done, then
\ui{Complete}, which disappears after two seconds. The \ui{LOG} of the window adds
\emph{Strain computation complete.} A project keeps the strain parameters, and the strain itself
when you save the results with it (Section~\ref{sec:projects}).

\ui{Cancel} appears only while the strain is being computed. It stops at the next frame and
keeps the previous strain result (\emph{Strain computation canceled.}). If the computation fails,
a \ui{Strain Computation Failed} dialog gives the reason.

\ui{Export Results} opens the export dialog with the display settings of this window
(Section~\ref{sec:export}). It is available as soon as there are displacement results; the strain
is exported once it is computed.

\subsubsection*{When the strain is out of date}

The strain on show was computed with the parameters of the last \ui{Compute Strain}. As soon as
you change a parameter, a warning line under the buttons reads \ui{Params changed -{}- click
Compute Strain}. Until you click it, the canvas, the \ui{Analysis} tab and exports keep using
the old strain. The line stays until the next \ui{Compute Strain}, even if you set the old values
back.

A new run of the analysis discards the strain, since it belongs to the previous displacements.
The strain buttons then show \emph{Click Compute Strain to show this field.} on the toolbar. A
strain computation still going on when a run finishes is discarded too (\emph{The displacement
results changed during the strain computation; its result was discarded.}).

\subsection{FIELD}
\label{sec:strain-fields}

\ui{FIELD} has two groups of buttons; one button is selected at a time.

Under \ui{DISPLACEMENT}, \ui{U}, \ui{V} and \ui{Magnitude} are the fields of the main window
(Section~\ref{sec:results}). \ui{Velocity} is the speed
$\sqrt{(\partial u/\partial t)^2 + (\partial v/\partial t)^2}$. Its time derivatives are central
differences of the displacement between the neighbouring computed frames, one-sided at the first
and last frame, so it assumes a constant interval between images. It is in pixels per frame, or
in the length unit per second with \ui{Use physical units} on. These four fields need no strain
computation.

Under \ui{STRAIN}, every field comes from the strain tensor of the chosen \ui{Strain type}, or
from $\mathbf{F}$ directly (Table~\ref{tab:strain-fields}). The strains are dimensionless; the two
angles are in degrees. The $y$ axis points down the image, so positive angles turn clockwise on
the screen.

\begin{table}[htbp]
\centering
\small
\caption{The strain fields. $\varepsilon$ stands for the strain tensor of the chosen \ui{Strain
type}.}
\label{tab:strain-fields}
\begin{tabularx}{\textwidth}{@{}l>{\raggedright\arraybackslash}X@{}}
\toprule
Button & Definition \\
\midrule
\ui{εxx}, \ui{εyy} & Normal strains along $x$ (right) and $y$ (down). \\
\ui{εxy} & Shear strain, the tensor component: half the engineering shear,
  $\varepsilon_{xy} = \gamma_{xy}/2$. \\
\ui{ε₁}, \ui{ε₂} & Major and minor principal strains,
  $\tfrac12(\varepsilon_{xx} + \varepsilon_{yy}) \pm
  \sqrt{\tfrac14(\varepsilon_{xx} - \varepsilon_{yy})^2 + \varepsilon_{xy}^2}$,
  with $\varepsilon_1 \geq \varepsilon_2$. \\
\ui{max shear} & Maximum in-plane shear strain, $(\varepsilon_1 - \varepsilon_2)/2$, also a
  tensor value (half the engineering shear). \\
\ui{von Mises} & Equivalent strain
  $\sqrt{\tfrac23(\varepsilon_1^2 + \varepsilon_2^2 + \varepsilon_3^2)}$ with
  $\varepsilon_3 = -(\varepsilon_1 + \varepsilon_2)$, which assumes an incompressible
  material. \\
\ui{ω rot} & Rigid rotation of the material, in degrees, from the polar decomposition of
  $\mathbf{F}$: $\operatorname{atan2}(F_{21} - F_{12},\, F_{11} + F_{22})$. \\
\ui{ε₁ angle} & Direction of $\varepsilon_1$ from $+x$, in degrees between $-90$ and 90:
  $\tfrac12\operatorname{atan2}(2\varepsilon_{xy},\, \varepsilon_{xx} - \varepsilon_{yy})$. \\
\ui{ΔA/A₀} & Relative change of area, $\det\mathbf{F} - 1$. \\
\bottomrule
\end{tabularx}
\end{table}

\subsection{VISUALIZATION}
\label{sec:strain-viz}

The controls work as in the main window (Section~\ref{sec:results}); this subsection lists what
differs.

\begin{description}[leftmargin=1.2cm,style=nextline]
  \item[Show on]
    \ui{Deformed frame} (the default) or \ui{Reference frame}. Strain is computed at the frame-1
    positions; the deformed view draws each value where its point is in the current frame, over
    that frame's image, as the main window does for displacement.
  \item[Background]
    \ui{Show background image}, with \ui{Hidden background} (\ui{White}, \ui{Black},
    \ui{Transparent}) for the hidden case. The window starts from the main window's
    \ui{Hidden background} when it is first opened.
  \item[Colormap, Range, Opacity]
    The same lists and rules as in the main window, remembered for each field. In \ui{Auto} the
    highest 2\,\% of the values, often a peak at the edge of a hole, take the end colour; set a
    \ui{Fixed} range to see the level of the peak. \ui{Min} and \ui{Max} hold six decimals here,
    enough for small strains. \ui{Opacity} is 70 by default.
  \item[Edges]
    \ui{Fill trimmed edges (display only)}, off by default, fills the trimmed band
    (Section~\ref{sec:edge-trim}) with the value of the nearest reliable point of the same
    region. Each region is filled from its own points, so the fill never crosses a crack or a
    hole. It applies to the view and to exported images and animations; exported data files
    always keep the band empty. The filled values are copies, not measurements: leave the option
    off when the edge itself matters.
  \item[Hide untracked regions]
    As in the main window; the window starts from the main window's choice when it is first
    opened, then keeps its own. On the \ui{Reference frame} view, unticking it fills every empty
    point inside the region of interest with the nearest value, the trimmed band included,
    whatever the \ui{Edges} option says. % CHECK: the untracked fill also fills the trimmed band (not region by region); intended? see report.
\end{description}

\ui{PHYSICAL UNITS} is the same group as in the main window and shares its setting. The window's
\ui{LOG} lists the strain computations; failures also open a dialog, since the \ui{LOG} is folded
by default.
